% Source: https://artofproblemsolving.com/community/c3h2559374_sum_of_logarithm_problem_i_made
% Stable identity: aops-topic-2559374-q01
Each term is $k\log\dfrac{k}{k+1}=k\log k-k\log(k+1)$. Re-indexing the second part with $m=k+1$,
\[\sum_{k=1}^{99}k\log k-\sum_{m=2}^{100}(m-1)\log m=1\cdot\log1+\sum_{m=2}^{99}\bigl(m-(m-1)\bigr)\log m-99\log100.\]
Since $\log1=0$ and $\log100=2$, the sum equals
\[\sum_{m=2}^{99}\log m-198=\boxed{\log(99!)-198}=\log\dfrac{99!}{10^{198}}.\]
Numerically, $\log(99!)\approx155.97$, so the sum is about $-42.03$. (In any base the sum is $\log\dfrac{99!}{100^{99}}$.)
